%% GENERATED by generate_papers.py -- edit the bank_*.py files, not this file.
\documentclass[11pt,a4paper]{article}
\usepackage[a4paper, top=20mm, bottom=20mm, left=18mm, right=18mm, headheight=14pt]{geometry}
\usepackage[T1]{fontenc}
\usepackage{lmodern}
\usepackage{amsmath, amssymb, mathtools}
\usepackage{braket}
\usepackage{array, tabularx, booktabs, multicol}
\usepackage{enumitem}
\usepackage[dvipsnames]{xcolor}
\usepackage[most]{tcolorbox}
\usepackage{needspace}
\usepackage{fancyhdr, lastpage}
\usepackage{pifont}
\usepackage{tikz}
\usepackage[hidelinks]{hyperref}
\newcolumntype{L}{>{\raggedright\arraybackslash}X}
\definecolor{brand}{HTML}{1F4E79}
\definecolor{okgreen}{HTML}{1B7F3B}
\setlength{\parindent}{0pt}
\setlength{\parskip}{3pt}
\newcommand{\optlabel}[1]{\textbf{(#1)}}
\newcommand{\correctlabel}[1]{\textcolor{okgreen}{\textbf{(#1)}\,\ding{51}}}
\newenvironment{qitem}[1]{\begin{list}{}{\setlength{\leftmargin}{3.3em}\setlength{\labelwidth}{3.1em}%
  \setlength{\labelsep}{0.2em}\setlength{\topsep}{3pt}\setlength{\parsep}{3pt}}\item[\textbf{Q.#1}\hfill]}{\end{list}}
\newcommand{\qsep}{\par\vspace{2pt}{\color{black!15}\hrule}\vspace{4pt}}
\newcommand{\qmeta}[2]{{\footnotesize\color{black!60}\textsc{#1} \;$\cdot$\; #2}\par}
\newtcolorbox{passagebox}[1]{enhanced, breakable, colback=brand!4, colframe=brand!70, fonttitle=\bfseries,
  title={#1}, left=6pt, right=6pt, top=4pt, bottom=4pt}
\newtcolorbox{sectionbanner}{enhanced, colback=brand, colframe=brand, coltext=white, fontupper=\large\bfseries,
  left=8pt, right=8pt, top=6pt, bottom=6pt, sharp corners}
\newtcolorbox{infobox}[1]{enhanced, breakable, colback=black!3, colframe=black!45, fonttitle=\bfseries,
  title={#1}, left=6pt, right=6pt, top=4pt, bottom=4pt}
\newtcolorbox{solbox}[1]{enhanced, breakable, colback=okgreen!3, colframe=okgreen!60!black, fonttitle=\bfseries,
  title={Answer: (#1) \quad Detailed solution}, left=6pt, right=6pt, top=3pt, bottom=4pt, before skip=4pt, after skip=6pt}
\newcommand{\bub}[1]{\tikz[baseline=-0.65ex]\node[draw, circle, inner sep=0pt, minimum size=4.6mm, font=\scriptsize]{#1};}
\allowdisplaybreaks

\begin{document}
\pagestyle{fancy}\fancyhf{}
\fancyhead[L]{\small\color{black!70}Grover's Algorithm --- Assessment}
\fancyhead[R]{\small\color{black!70}\textbf{Question Paper}}
\fancyfoot[C]{\small Page \thepage\ of \pageref*{LastPage}}
\renewcommand{\headrulewidth}{0.4pt}

\thispagestyle{empty}
\begin{center}
{\color{brand}\rule{\linewidth}{1.2pt}}\\[6pt]
{\Huge\bfseries\color{brand} Grover's Algorithm}\\[4pt]
{\Large Assessment of quantum search}\\[4pt]
{\large\bfseries QUESTION PAPER}\\[2pt]
{\color{brand}\rule{\linewidth}{1.2pt}}\\[1pt]
\end{center}
\vspace{-2pt}
\begin{tabularx}{\linewidth}{@{}lXlr@{}}
\textbf{Time allowed:} & 3 hours & \textbf{Maximum marks:} & 240\\
\textbf{Questions:} & 80 (4 sections $\times$ 20) & \textbf{Answer key:} & separate Answer Book
\end{tabularx}
\vspace{6pt}

\textbf{Name of the candidate:} \hrulefill\quad \textbf{Roll No.:} \rule{35mm}{0.4pt}

\vspace{6pt}
\begin{infobox}{General instructions}
\begin{enumerate}[leftmargin=1.4em, itemsep=1pt, topsep=1pt]
\item This question paper has \textbf{four sections} of \textbf{20 questions} each:
\textbf{A}~Mathematics for Grover, \textbf{B}~Quantum Mechanics, \textbf{C}~Grover Basics, \textbf{D}~Advanced Grover.
\item Every question is of the \textbf{single-correct} type: four options, exactly one correct. Question formats used:
single-correct MCQ, matching lists (List-I/List-II), Assertion--Reason, Statement~I/II and paragraph-based questions.
\item Mark your answers on the \textbf{Response Sheet} at the end of this paper by darkening ONE bubble per question.
\item Calculators are \textbf{not} allowed. Rough work may be done on separate sheets.
\item Unless a question says otherwise, use the \textbf{conventions} printed below; they apply to every section.
\item The detailed solutions are in the separate \emph{Answer Book}; do not consult it before finishing the test.
\end{enumerate}
\end{infobox}
\begin{infobox}{Marking scheme (every question)}
Each question has \textbf{four} options (A), (B), (C) and (D). \textbf{ONLY ONE} of these four options is correct.\\[2pt]
\begin{tabular}{@{}ll@{}}
\emph{Full marks}: & $+3$ if ONLY the correct option is chosen;\\
\emph{Zero marks}: & $\phantom{+}0$ if none of the options is chosen (question unanswered);\\
\emph{Negative marks}: & $-1$ in all other cases (wrong option, or more than one option chosen).
\end{tabular}
\end{infobox}

\begin{infobox}{Conventions and notation used throughout the paper}
\small
\begin{itemize}[leftmargin=1.2em, itemsep=1pt, topsep=2pt]
\item $n$ qubits, $N = 2^n$ items, $M$ marked items ($M = 1$ unless stated otherwise). A basis state is written $\ket{x} = \ket{x_1x_2\dots x_n}$:
      \textbf{qubit 1 is the most significant bit} and is written first. (The accompanying notebook writes \texttt{N\_QUBITS} for $n$ and $M = 2^{N}$ for the number of items.)
\item $\ket{s} = H^{\otimes n}\ket{0\dots0} = \frac{1}{\sqrt N}\sum_x\ket{x}$. The oracle is $O\ket{x} = (-1)^{f(x)}\ket{x}$ with $f(x) = 1$ exactly on the marked items; the diffusion is $D = 2\ket{s}\!\bra{s} - I$; one iteration is $G = DO$ (oracle first).
\item $\sin\theta = \sqrt{M/N}$ with $0 < \theta \le \frac\pi2$; $\;\ket{m}$ and $\ket{r}$ are the uniform superpositions of the marked and the unmarked items. Then
      $\displaystyle G^k\ket{s} = \sin\big((2k+1)\theta\big)\ket{m} + \cos\big((2k+1)\theta\big)\ket{r}$, so $P_k = \sin^2\big((2k+1)\theta\big)$ and
      $k_{\rm opt} = \operatorname{round}\big(\frac{\pi}{4\theta} - \frac12\big)$.
\item \textbf{Textbook circuit:} $H^{\otimes n}$, then $k$ times: [$X$ on each qubit where $x^\ast$ has a $0$; a multi-controlled $Z$ on all $n$ qubits; the same $X$ gates] followed by [$H^{\otimes n}$; $X^{\otimes n}$; a multi-controlled $Z$; $X^{\otimes n}$; $H^{\otimes n}$]; finally measure all qubits.
\item $H = \frac{1}{\sqrt2}\begin{psmallmatrix}1&1\\1&-1\end{psmallmatrix}$; $X, Y, Z$ are the Pauli matrices; $\ket{\pm} = \frac{1}{\sqrt2}(\ket0 \pm \ket1)$. Two-qubit matrices use the basis order $\ket{00}, \ket{01}, \ket{10}, \ket{11}$.
\item All measurements are in the computational basis unless stated otherwise; $\log$ means $\log_2$. A numerical value given in a question may be used without a calculator.
\end{itemize}
\end{infobox}
\newpage

\begin{sectionbanner}SECTION A \enspace$\cdot$\enspace Mathematics for Grover \hfill Q.1--20\end{sectionbanner}

{\small \textbf{Syllabus:} trigonometric identities, reflections and rotations in the plane, eigenvalues, inversion about the mean, Boolean counting, classical search probability and square-root scaling.\\ \textbf{Formats:} Single-correct MCQ (Q.1--14); Matching lists (Q.15--16); Assertion--Reason (Q.17); Statement I / II (Q.18); Paragraph based (Q.19--20). Each question: $+3$ / $0$ / $-1$.}
\par\medskip

\needspace{7\baselineskip}
\begin{qitem}{1}
If $\sin\theta = \tfrac14$, the value of $\sin^2 3\theta$ is
\par\smallskip\noindent\begin{tabularx}{\linewidth}{@{}*{4}{L}@{}}
\optlabel{A}\,$\dfrac{1}{16}$ & \optlabel{B}\,$\dfrac{121}{256}$ & \optlabel{C}\,$\dfrac{9}{16}$ & \optlabel{D}\,$\dfrac{11}{16}$
\end{tabularx}
\end{qitem}
\qsep

\needspace{7\baselineskip}
\begin{qitem}{2}
Let $u = \big(\tfrac35, \tfrac45\big)^{T}$. The matrix $2uu^{T} - I$ equals
\par\smallskip\noindent\begin{tabularx}{\linewidth}{@{}*{4}{L}@{}}
\optlabel{A}\,$\dfrac{1}{25}\begin{pmatrix}18 & 24\\ 24 & 32\end{pmatrix}$ & \optlabel{B}\,$\dfrac{1}{25}\begin{pmatrix}-7 & 12\\ 12 & 7\end{pmatrix}$ & \optlabel{C}\,$\dfrac{1}{25}\begin{pmatrix}-7 & 24\\ 24 & 7\end{pmatrix}$ & \optlabel{D}\,$\dfrac{1}{25}\begin{pmatrix}7 & -24\\ -24 & -7\end{pmatrix}$
\end{tabularx}
\end{qitem}
\qsep

\needspace{7\baselineskip}
\begin{qitem}{3}
A vector in the plane is reflected first in the line through the origin at angle $10^\circ$ to the $x$-axis, and then in the
line through the origin at angle $35^\circ$. The combined map is
\par\smallskip\noindent\begin{tabularx}{\linewidth}{@{}*{2}{L}@{}}
\optlabel{A}\,a reflection in the line at angle $45^\circ$ & \optlabel{B}\,a rotation by $+25^\circ$ \\[3pt]
\optlabel{C}\,a rotation by $-50^\circ$ & \optlabel{D}\,a rotation by $+50^\circ$
\end{tabularx}
\end{qitem}
\qsep

\needspace{7\baselineskip}
\begin{qitem}{4}
Let $\sin\theta = \tfrac14$ with $0 < \theta < \tfrac\pi2$, and let
$G = \begin{psmallmatrix}\cos 2\theta & -\sin 2\theta\\ \sin 2\theta & \cos 2\theta\end{psmallmatrix}$. The eigenvalues of $G$ are
\par\smallskip\noindent\begin{tabularx}{\linewidth}{@{}*{4}{L}@{}}
\optlabel{A}\,$\dfrac{7 \pm \sqrt{15}}{8}$ & \optlabel{B}\,$\dfrac{1 \pm i\sqrt{15}}{4}$ & \optlabel{C}\,$\pm 1$ & \optlabel{D}\,$\dfrac{7 \pm i\sqrt{15}}{8}$
\end{tabularx}
\end{qitem}
\qsep

\needspace{7\baselineskip}
\begin{qitem}{5}
Let $R = \begin{psmallmatrix}\sqrt3/2 & -1/2\\ 1/2 & \sqrt3/2\end{psmallmatrix}$. The smallest positive integer $k$ with $R^k = -I$ is
\par\smallskip\noindent\begin{tabularx}{\linewidth}{@{}*{4}{L}@{}}
\optlabel{A}\,$6$ & \optlabel{B}\,$4$ & \optlabel{C}\,$3$ & \optlabel{D}\,$12$
\end{tabularx}
\end{qitem}
\qsep

\needspace{7\baselineskip}
\begin{qitem}{6}
Let $\sin\theta = \tfrac{1}{10}$, $0 < \theta < \tfrac\pi2$; numerically $\theta = 0.10017$. Among the integers $0 \le k \le 10$, the one that makes
$\sin^2\big((2k+1)\theta\big)$ largest is
\par\smallskip\noindent\begin{tabularx}{\linewidth}{@{}*{4}{L}@{}}
\optlabel{A}\,$8$ & \optlabel{B}\,$10$ & \optlabel{C}\,$5$ & \optlabel{D}\,$7$
\end{tabularx}
\end{qitem}
\qsep

\needspace{7\baselineskip}
\begin{qitem}{7}
The vector $v \in \mathbb{R}^{8}$ has seven entries equal to $+1$ and one entry equal to $-1$. Each entry $v_j$ is replaced by $2\bar v - v_j$,
where $\bar v$ is the mean of the entries. The new value of the entry that was $-1$ is
\par\smallskip\noindent\begin{tabularx}{\linewidth}{@{}*{4}{L}@{}}
\optlabel{A}\,$\dfrac52$ & \optlabel{B}\,$\dfrac32$ & \optlabel{C}\,$\dfrac74$ & \optlabel{D}\,$\dfrac12$
\end{tabularx}
\end{qitem}
\qsep

\needspace{7\baselineskip}
\begin{qitem}{8}
A box holds $31$ cards of which exactly $3$ are marked. Cards are drawn one at a time at random without replacement until the first
marked card appears. The expected number of draws is
\par\smallskip\noindent\begin{tabularx}{\linewidth}{@{}*{4}{L}@{}}
\optlabel{A}\,$\dfrac{31}{4}$ & \optlabel{B}\,$\dfrac{31}{3}$ & \optlabel{C}\,$8$ & \optlabel{D}\,$16$
\end{tabularx}
\end{qitem}
\qsep

\needspace{7\baselineskip}
\begin{qitem}{9}
One run of a procedure succeeds with probability $\tfrac14$, independently of the other runs. The probability of at least one success in
three runs is
\par\smallskip\noindent\begin{tabularx}{\linewidth}{@{}*{4}{L}@{}}
\optlabel{A}\,$\dfrac{37}{64}$ & \optlabel{B}\,$\dfrac{27}{64}$ & \optlabel{C}\,$\dfrac34$ & \optlabel{D}\,$\dfrac{1}{64}$
\end{tabularx}
\end{qitem}
\qsep

\needspace{7\baselineskip}
\begin{qitem}{10}
Let $f(x_1, x_2, x_3, x_4) = (x_1 \lor x_2) \land (x_3 \oplus x_4)$ on $\{0, 1\}^4$. The fraction of inputs with $f = 1$ is
\par\smallskip\noindent\begin{tabularx}{\linewidth}{@{}*{4}{L}@{}}
\optlabel{A}\,$\dfrac12$ & \optlabel{B}\,$\dfrac38$ & \optlabel{C}\,$\dfrac14$ & \optlabel{D}\,$\dfrac58$
\end{tabularx}
\end{qitem}
\qsep

\needspace{7\baselineskip}
\begin{qitem}{11}
A Boolean function $f$ on $4$ bits takes the value $1$ on exactly $3$ inputs. The value of $\displaystyle\sum_{x\in\{0,1\}^4}(-1)^{f(x)}$ is
\par\smallskip\noindent\begin{tabularx}{\linewidth}{@{}*{4}{L}@{}}
\optlabel{A}\,$10$ & \optlabel{B}\,$7$ & \optlabel{C}\,$6$ & \optlabel{D}\,$13$
\end{tabularx}
\end{qitem}
\qsep

\needspace{7\baselineskip}
\begin{qitem}{12}
If $\sin^2\theta = \tfrac{1}{16}$, then $\cos 4\theta$ equals
\par\smallskip\noindent\begin{tabularx}{\linewidth}{@{}*{4}{L}@{}}
\optlabel{A}\,$\dfrac12$ & \optlabel{B}\,$\dfrac{17}{32}$ & \optlabel{C}\,$\dfrac34$ & \optlabel{D}\,$\dfrac{49}{64}$
\end{tabularx}
\end{qitem}
\qsep

\needspace{7\baselineskip}
\begin{qitem}{13}
Let $g(n) = \tfrac\pi4\sqrt{2^n}$. Then $g(n + 6)/g(n)$ equals
\par\smallskip\noindent\begin{tabularx}{\linewidth}{@{}*{4}{L}@{}}
\optlabel{A}\,$6$ & \optlabel{B}\,$8$ & \optlabel{C}\,$\sqrt6$ & \optlabel{D}\,$64$
\end{tabularx}
\end{qitem}
\qsep

\needspace{7\baselineskip}
\begin{qitem}{14}
Let $\theta = \pi/10$. The smallest positive integer $p$ such that $\sin^2\big((2k + 2p + 1)\theta\big) = \sin^2\big((2k+1)\theta\big)$ for every integer $k$ is
\par\smallskip\noindent\begin{tabularx}{\linewidth}{@{}*{4}{L}@{}}
\optlabel{A}\,$10$ & \optlabel{B}\,$20$ & \optlabel{C}\,$5$ & \optlabel{D}\,$2$
\end{tabularx}
\end{qitem}
\qsep

\needspace{7\baselineskip}
\begin{qitem}{15}
Match the value of $\sin^2\theta$ (List-I) with the value of $\sin^2 3\theta$ (List-II).
\par\smallskip\noindent{\renewcommand{\arraystretch}{1.35}%
\begin{tabularx}{\linewidth}{@{}c>{\raggedright\arraybackslash}p{0.33\linewidth}@{\hspace{1.2em}}cL@{}}
\toprule
\multicolumn{2}{@{}l}{\textbf{List-I}} & \multicolumn{2}{l@{}}{\textbf{List-II}} \\ \midrule
\textbf{(P)} & $\sin^2\theta = \frac18$ & \textbf{(1)} & $1$ \\
\textbf{(Q)} & $\sin^2\theta = \frac14$ & \textbf{(2)} & $0$ \\
\textbf{(R)} & $\sin^2\theta = \frac12$ & \textbf{(3)} & $\frac{25}{32}$ \\
\textbf{(S)} & $\sin^2\theta = \frac34$ & \textbf{(4)} & $\frac12$ \\
\bottomrule
\end{tabularx}}
\par\smallskip\noindent\begin{tabularx}{\linewidth}{@{}*{2}{L}@{}}
\optlabel{A}\,P$\to$4, Q$\to$1, R$\to$3, S$\to$2 & \optlabel{B}\,P$\to$3, Q$\to$1, R$\to$2, S$\to$4 \\[3pt]
\optlabel{C}\,P$\to$3, Q$\to$1, R$\to$4, S$\to$2 & \optlabel{D}\,P$\to$1, Q$\to$3, R$\to$4, S$\to$2
\end{tabularx}
\end{qitem}
\qsep

\needspace{7\baselineskip}
\begin{qitem}{16}
Let $\ket{u}$ be a unit vector in $\mathbb{C}^{N}$, with $N \ge 3$. Match each matrix in List-I with its trace in List-II.
\par\smallskip\noindent{\renewcommand{\arraystretch}{1.35}%
\begin{tabularx}{\linewidth}{@{}c>{\raggedright\arraybackslash}p{0.33\linewidth}@{\hspace{1.2em}}cL@{}}
\toprule
\multicolumn{2}{@{}l}{\textbf{List-I}} & \multicolumn{2}{l@{}}{\textbf{List-II}} \\ \midrule
\textbf{(P)} & $2\ket{u}\!\bra{u} - I$ & \textbf{(1)} & $N$ \\
\textbf{(Q)} & $I - 2\ket{u}\!\bra{u}$ & \textbf{(2)} & $N - 2$ \\
\textbf{(R)} & $\ket{u}\!\bra{u}$ & \textbf{(3)} & $1$ \\
\textbf{(S)} & $I$ & \textbf{(4)} & $2 - N$ \\
\bottomrule
\end{tabularx}}
\par\smallskip\noindent\begin{tabularx}{\linewidth}{@{}*{2}{L}@{}}
\optlabel{A}\,P$\to$3, Q$\to$2, R$\to$4, S$\to$1 & \optlabel{B}\,P$\to$4, Q$\to$2, R$\to$1, S$\to$3 \\[3pt]
\optlabel{C}\,P$\to$2, Q$\to$4, R$\to$3, S$\to$1 & \optlabel{D}\,P$\to$4, Q$\to$2, R$\to$3, S$\to$1
\end{tabularx}
\end{qitem}
\qsep

\needspace{7\baselineskip}
\begin{qitem}{17}
\textbf{Assertion (A):} Rounding $\frac\pi4\sqrt N$ to the nearest integer does not always give the integer $k$ for which
$(2k+1)\theta$ is closest to $\frac\pi2$, where $\sin\theta = 1/\sqrt N$.\\[1mm]
\textbf{Reason (R):} For $0 < x < 1$ we have $\arcsin x < x$.
\begin{itemize}[leftmargin=2.6em, labelsep=0.4em, itemsep=2pt, topsep=3pt]
\item[\optlabel{A}] Both Assertion (A) and Reason (R) are true, and R is the correct explanation of A
\item[\optlabel{B}] Both A and R are true, but R is \emph{not} the correct explanation of A
\item[\optlabel{C}] A is true, but R is false
\item[\optlabel{D}] A is false, but R is true
\end{itemize}
\end{qitem}
\qsep

\needspace{7\baselineskip}
\begin{qitem}{18}
Let $\ket{u}$ be a unit vector in $\mathbb{C}^{N}$, with $N \ge 2$, and let $D = 2\ket{u}\!\bra{u} - I$.\\[1mm]
\textbf{Statement I:} $D^2 = I$.\\[1mm]
\textbf{Statement II:} $D$ has the eigenvalue $+1$ with multiplicity $1$ and the eigenvalue $-1$ with multiplicity $N - 1$.
\begin{itemize}[leftmargin=2.6em, labelsep=0.4em, itemsep=2pt, topsep=3pt]
\item[\optlabel{A}] Both Statement I and Statement II are true
\item[\optlabel{B}] Both Statement I and Statement II are false
\item[\optlabel{C}] Statement I is true, but Statement II is false
\item[\optlabel{D}] Statement I is false, but Statement II is true
\end{itemize}
\end{qitem}
\qsep

\needspace{20\baselineskip}
\begin{passagebox}{Paragraph for Questions 19 and 20}
Define $p(k) = \sin^2\big((2k+1)\theta\big)$ for integers $k \ge 0$, with $\theta = \pi/18 = 10^\circ$. (This is the
shape of the Grover success curve, but the two questions below need only trigonometry.) Note that $(2k+1)\theta$ takes the
values $10^\circ, 30^\circ, 50^\circ, \dots$
\end{passagebox}
\needspace{7\baselineskip}
\begin{qitem}{19}
The smallest $k$ with $p(k) = 1$, and the smallest positive period of the sequence $p(k)$, are
\begin{itemize}[leftmargin=2.6em, labelsep=0.4em, itemsep=2pt, topsep=3pt]
\item[\optlabel{A}] first maximum at $k = 4$; smallest period $18$
\item[\optlabel{B}] first maximum at $k = 5$; smallest period $9$
\item[\optlabel{C}] first maximum at $k = 9$; smallest period $18$
\item[\optlabel{D}] first maximum at $k = 4$; smallest period $9$
\end{itemize}
\end{qitem}
\qsep

\needspace{7\baselineskip}
\begin{qitem}{20}
The number of integers $k$ with $0 \le k \le 17$ for which $p(k) \ge \tfrac34$ is
\par\smallskip\noindent\begin{tabularx}{\linewidth}{@{}*{4}{L}@{}}
\optlabel{A}\,$3$ & \optlabel{B}\,$6$ & \optlabel{C}\,$9$ & \optlabel{D}\,$4$
\end{tabularx}
\end{qitem}
\qsep

\newpage

\begin{sectionbanner}SECTION B \enspace$\cdot$\enspace Quantum Mechanics \hfill Q.21--40\end{sectionbanner}

{\small \textbf{Syllabus:} superposition, Born rule, global and relative phases, Hadamard layers, reflections, bit and phase oracles, phase kickback, multi-controlled gates, entanglement and partial measurement.\\ \textbf{Formats:} Single-correct MCQ (Q.21--34); Matching lists (Q.35--36); Assertion--Reason (Q.37); Statement I / II (Q.38); Paragraph based (Q.39--40). Each question: $+3$ / $0$ / $-1$.}
\par\medskip

\needspace{7\baselineskip}
\begin{qitem}{21}
The state $H^{\otimes3}\ket{000}$ is measured in the computational basis. The probability that qubit 1 gives $1$ and qubit 3 gives $0$ is
\par\smallskip\noindent\begin{tabularx}{\linewidth}{@{}*{4}{L}@{}}
\optlabel{A}\,$\dfrac18$ & \optlabel{B}\,$\dfrac14$ & \optlabel{C}\,$\dfrac{1}{\sqrt8}$ & \optlabel{D}\,$\dfrac12$
\end{tabularx}
\end{qitem}
\qsep

\needspace{7\baselineskip}
\begin{qitem}{22}
The state $\frac12\big(\ket{00} + \ket{01} + \ket{10} - \ket{11}\big)$ is measured in the computational basis. The probability of $\ket{11}$ is
\par\smallskip\noindent\begin{tabularx}{\linewidth}{@{}*{4}{L}@{}}
\optlabel{A}\,$\dfrac14$ & \optlabel{B}\,$0$ & \optlabel{C}\,$\dfrac12$ & \optlabel{D}\,$-\dfrac14$
\end{tabularx}
\end{qitem}
\qsep

\needspace{7\baselineskip}
\begin{qitem}{23}
$H \otimes H$ is applied to $\frac12\big(\ket{00} + \ket{01} + \ket{10} - \ket{11}\big)$. The amplitude of $\ket{11}$ in the result is
\par\smallskip\noindent\begin{tabularx}{\linewidth}{@{}*{4}{L}@{}}
\optlabel{A}\,$\dfrac12$ & \optlabel{B}\,$0$ & \optlabel{C}\,$-\dfrac12$ & \optlabel{D}\,$1$
\end{tabularx}
\end{qitem}
\qsep

\needspace{7\baselineskip}
\begin{qitem}{24}
Let $U_f\ket{x}\ket{y} = \ket{x}\ket{y \oplus f(x)}$ and $\ket{-} = \frac{1}{\sqrt2}(\ket0 - \ket1)$. Then $U_f\ket{x}\ket{-}$ equals
\par\smallskip\noindent\begin{tabularx}{\linewidth}{@{}*{4}{L}@{}}
\optlabel{A}\,$\ket{x}\ket{f(x)}$ & \optlabel{B}\,$\ket{x}\ket{-}$ & \optlabel{C}\,$(-1)^{f(x)}\ket{x}\ket{-}$ & \optlabel{D}\,$\ket{f(x)}\ket{-}$
\end{tabularx}
\end{qitem}
\qsep

\needspace{7\baselineskip}
\begin{qitem}{25}
With $U_f$ as in the previous question and $\ket{+} = \frac{1}{\sqrt2}(\ket0 + \ket1)$, the state $U_f\ket{x}\ket{+}$ equals
\par\smallskip\noindent\begin{tabularx}{\linewidth}{@{}*{4}{L}@{}}
\optlabel{A}\,$\ket{x}\ket{-}$ & \optlabel{B}\,$(-1)^{f(x)}\ket{x}\ket{-}$ & \optlabel{C}\,$(-1)^{f(x)}\ket{x}\ket{+}$ & \optlabel{D}\,$\ket{x}\ket{+}$
\end{tabularx}
\end{qitem}
\qsep

\needspace{7\baselineskip}
\begin{qitem}{26}
Let CCZ be the three-qubit gate that multiplies only $\ket{111}$ by $-1$. Which basis state has its sign flipped by $X_2\,\mathrm{CCZ}\,X_2$, where $X_2$ is $X$ on qubit 2?
\par\smallskip\noindent\begin{tabularx}{\linewidth}{@{}*{4}{L}@{}}
\optlabel{A}\,$\ket{111}$ & \optlabel{B}\,$\ket{010}$ & \optlabel{C}\,$\ket{101}$ & \optlabel{D}\,$\ket{000}$
\end{tabularx}
\end{qitem}
\qsep

\needspace{7\baselineskip}
\begin{qitem}{27}
Let CCX be the Toffoli gate with target qubit 3. The product $(I \otimes I \otimes H)\,\mathrm{CCX}\,(I \otimes I \otimes H)$ equals
\par\smallskip\noindent\begin{tabularx}{\linewidth}{@{}*{4}{L}@{}}
\optlabel{A}\,CCX & \optlabel{B}\,$I \otimes I \otimes I$ & \optlabel{C}\,$Z \otimes Z \otimes Z$ & \optlabel{D}\,CCZ
\end{tabularx}
\end{qitem}
\qsep

\needspace{7\baselineskip}
\begin{qitem}{28}
Which pair of states are physically \emph{indistinguishable} (they give the same statistics for every possible measurement)?
\begin{itemize}[leftmargin=2.6em, labelsep=0.4em, itemsep=2pt, topsep=3pt]
\item[\optlabel{A}] $\frac{1}{\sqrt2}(\ket0 + \ket1)$ and $-\frac{1}{\sqrt2}(\ket0 + \ket1)$
\item[\optlabel{B}] $\frac{1}{\sqrt2}(\ket{00} + \ket{11})$ and $\frac{1}{\sqrt2}(\ket{00} - \ket{11})$
\item[\optlabel{C}] $\frac{1}{\sqrt2}(\ket0 + \ket1)$ and $\frac{1}{\sqrt2}(\ket0 - \ket1)$
\item[\optlabel{D}] $\frac{1}{\sqrt2}(\ket0 + \ket1)$ and $\frac{1}{\sqrt2}(\ket0 + i\ket1)$
\end{itemize}
\end{qitem}
\qsep

\needspace{7\baselineskip}
\begin{qitem}{29}
Let $\ket\psi = \frac{1}{\sqrt{10}}\big(\ket{00} + 2\ket{01} - 2\ket{10} + i\ket{11}\big)$. A measurement asks only whether the two qubits have \emph{different} values
(outcome ``different'' on the span of $\ket{01}, \ket{10}$). The probability of ``different'' and the state afterwards are
\par\smallskip\noindent\begin{tabularx}{\linewidth}{@{}*{2}{L}@{}}
\optlabel{A}\,$\frac45$ and $\ket{01}$ & \optlabel{B}\,$\frac45$ and $\frac{1}{\sqrt2}(\ket{01} - \ket{10})$ \\[3pt]
\optlabel{C}\,$\frac45$ and $\frac{1}{\sqrt{10}}(2\ket{01} - 2\ket{10})$ & \optlabel{D}\,$\frac25$ and $\frac{1}{\sqrt2}(\ket{01} - \ket{10})$
\end{tabularx}
\end{qitem}
\qsep

\needspace{7\baselineskip}
\begin{qitem}{30}
Let $\ket{u}$ be a unit vector in a space of dimension $N \ge 2$. Which of these operators is \emph{not} unitary?
\par\smallskip\noindent\begin{tabularx}{\linewidth}{@{}*{2}{L}@{}}
\optlabel{A}\,$H^{\otimes n}$ (for $N = 2^n$) & \optlabel{B}\,$\ket{u}\!\bra{u}$ \\[3pt]
\optlabel{C}\,$I - 2\ket{u}\!\bra{u}$ & \optlabel{D}\,$2\ket{u}\!\bra{u} - I$
\end{tabularx}
\end{qitem}
\qsep

\needspace{7\baselineskip}
\begin{qitem}{31}
For $N = 16$ let $\ket{s}$ be the uniform superposition of all $16$ basis states and $\ket{r}$ the uniform superposition of $15$ of them (all except one). The value of $\braket{r|s}$ is
\par\smallskip\noindent\begin{tabularx}{\linewidth}{@{}*{4}{L}@{}}
\optlabel{A}\,$\dfrac{\sqrt{15}}{4}$ & \optlabel{B}\,$\dfrac{15}{16}$ & \optlabel{C}\,$\dfrac14$ & \optlabel{D}\,$\dfrac{1}{\sqrt{15}}$
\end{tabularx}
\end{qitem}
\qsep

\needspace{7\baselineskip}
\begin{qitem}{32}
A state vector of $n = 33$ qubits is stored in double precision ($16$ bytes per complex amplitude). The memory needed is
\par\smallskip\noindent\begin{tabularx}{\linewidth}{@{}*{4}{L}@{}}
\optlabel{A}\,$256$ GiB & \optlabel{B}\,$16$ GiB & \optlabel{C}\,$64$ GiB & \optlabel{D}\,$128$ GiB
\end{tabularx}
\end{qitem}
\qsep

\needspace{7\baselineskip}
\begin{qitem}{33}
For a single qubit ($n = 1$), the operator $2\ket{+}\!\bra{+} - I$ equals
\par\smallskip\noindent\begin{tabularx}{\linewidth}{@{}*{4}{L}@{}}
\optlabel{A}\,$X$ & \optlabel{B}\,$H$ & \optlabel{C}\,$Z$ & \optlabel{D}\,$-I$
\end{tabularx}
\end{qitem}
\qsep

\needspace{7\baselineskip}
\begin{qitem}{34}
For $n$ qubits let $\ket{s} = H^{\otimes n}\ket{0\dots0}$ and $D = 2\ket{s}\!\bra{s} - I$. If the oracle marks nothing (it acts as the identity), then after $k$ iterations the state is
\par\smallskip\noindent\begin{tabularx}{\linewidth}{@{}*{2}{L}@{}}
\optlabel{A}\,$(-1)^k\ket{s}$ & \optlabel{B}\,$\ket{0\dots0}$ for odd $k$ \\[3pt]
\optlabel{C}\,a state whose amplitudes depend on $k$ & \optlabel{D}\,$\ket{s}$ for every $k$
\end{tabularx}
\end{qitem}
\qsep

\needspace{7\baselineskip}
\begin{qitem}{35}
Match each expression in List-I with its value in List-II.
\par\smallskip\noindent{\renewcommand{\arraystretch}{1.35}%
\begin{tabularx}{\linewidth}{@{}c>{\raggedright\arraybackslash}p{0.33\linewidth}@{\hspace{1.2em}}cL@{}}
\toprule
\multicolumn{2}{@{}l}{\textbf{List-I}} & \multicolumn{2}{l@{}}{\textbf{List-II}} \\ \midrule
\textbf{(P)} & $H\ket1$ & \textbf{(1)} & $-\ket{-}$ \\
\textbf{(Q)} & $Z\ket{-}$ & \textbf{(2)} & $\ket0$ \\
\textbf{(R)} & $X\ket{-}$ & \textbf{(3)} & $\ket{+}$ \\
\textbf{(S)} & $H\ket{+}$ & \textbf{(4)} & $\ket{-}$ \\
\bottomrule
\end{tabularx}}
\par\smallskip\noindent\begin{tabularx}{\linewidth}{@{}*{2}{L}@{}}
\optlabel{A}\,P$\to$4, Q$\to$1, R$\to$3, S$\to$2 & \optlabel{B}\,P$\to$4, Q$\to$3, R$\to$2, S$\to$1 \\[3pt]
\optlabel{C}\,P$\to$4, Q$\to$3, R$\to$1, S$\to$2 & \optlabel{D}\,P$\to$3, Q$\to$4, R$\to$1, S$\to$2
\end{tabularx}
\end{qitem}
\qsep

\needspace{7\baselineskip}
\begin{qitem}{36}
Match each two-qubit state in List-I with its description in List-II.
\par\smallskip\noindent{\renewcommand{\arraystretch}{1.35}%
\begin{tabularx}{\linewidth}{@{}c>{\raggedright\arraybackslash}p{0.33\linewidth}@{\hspace{1.2em}}cL@{}}
\toprule
\multicolumn{2}{@{}l}{\textbf{List-I}} & \multicolumn{2}{l@{}}{\textbf{List-II}} \\ \midrule
\textbf{(P)} & $\frac12(\ket{00} + \ket{01} + \ket{10} + \ket{11})$ & \textbf{(1)} & entangled, with four amplitudes of equal size \\
\textbf{(Q)} & $\frac12(\ket{00} - \ket{01} - \ket{10} + \ket{11})$ & \textbf{(2)} & the product $\ket{-}\ket{-}$ \\
\textbf{(R)} & $\frac12(\ket{00} + \ket{01} + \ket{10} - \ket{11})$ & \textbf{(3)} & entangled, with only two non-zero amplitudes \\
\textbf{(S)} & $\frac{1}{\sqrt2}(\ket{00} + \ket{11})$ & \textbf{(4)} & the product $\ket{+}\ket{+}$ \\
\bottomrule
\end{tabularx}}
\par\smallskip\noindent\begin{tabularx}{\linewidth}{@{}*{2}{L}@{}}
\optlabel{A}\,P$\to$4, Q$\to$2, R$\to$3, S$\to$1 & \optlabel{B}\,P$\to$4, Q$\to$2, R$\to$1, S$\to$3 \\[3pt]
\optlabel{C}\,P$\to$4, Q$\to$1, R$\to$2, S$\to$3 & \optlabel{D}\,P$\to$2, Q$\to$4, R$\to$1, S$\to$3
\end{tabularx}
\end{qitem}
\qsep

\needspace{7\baselineskip}
\begin{qitem}{37}
\textbf{Assertion (A):} Measuring in the computational basis immediately after a phase oracle acts on $\ket{s}$ gives a uniformly random string; the oracle step alone reveals nothing.\\[1mm]
\textbf{Reason (R):} The oracle multiplies each amplitude by $\pm1$, and the Born probabilities depend only on the magnitudes of the amplitudes.
\begin{itemize}[leftmargin=2.6em, labelsep=0.4em, itemsep=2pt, topsep=3pt]
\item[\optlabel{A}] Both Assertion (A) and Reason (R) are true, and R is the correct explanation of A
\item[\optlabel{B}] Both A and R are true, but R is \emph{not} the correct explanation of A
\item[\optlabel{C}] A is true, but R is false
\item[\optlabel{D}] A is false, but R is true
\end{itemize}
\end{qitem}
\qsep

\needspace{7\baselineskip}
\begin{qitem}{38}
\textbf{Statement I:} The states $\ket\psi$ and $-\ket\psi$ give identical statistics for every measurement.\\[1mm]
\textbf{Statement II:} Multiplying only \emph{one} basis-state amplitude of a superposition by $-1$ can never change the statistics of any measurement.
\begin{itemize}[leftmargin=2.6em, labelsep=0.4em, itemsep=2pt, topsep=3pt]
\item[\optlabel{A}] Both Statement I and Statement II are true
\item[\optlabel{B}] Both Statement I and Statement II are false
\item[\optlabel{C}] Statement I is true, but Statement II is false
\item[\optlabel{D}] Statement I is false, but Statement II is true
\end{itemize}
\end{qitem}
\qsep

\needspace{20\baselineskip}
\begin{passagebox}{Paragraph for Questions 39 and 40}
Three qubits start in $\ket{0}\ket{0}\ket{1}$ (input qubits 1 and 2, ancilla qubit 3). A Hadamard gate is applied to all three. Then the
Toffoli gate is applied with qubits 1 and 2 as controls and qubit 3 as target: it implements
$U_f\ket{x}\ket{y} = \ket{x}\ket{y \oplus f(x)}$ with $f(x_1x_2) = x_1 \land x_2$.
\end{passagebox}
\needspace{7\baselineskip}
\begin{qitem}{39}
The state of the three qubits after the Toffoli gate is
\par\smallskip\noindent\begin{tabularx}{\linewidth}{@{}*{2}{L}@{}}
\optlabel{A}\,$\frac12\big(\ket{00} - \ket{01} - \ket{10} - \ket{11}\big)\ket{-}$ & \optlabel{B}\,$\frac12\big(\ket{00} + \ket{01} + \ket{10} - \ket{11}\big)\ket{+}$ \\[3pt]
\optlabel{C}\,$\frac12\big(\ket{00} + \ket{01} + \ket{10} + \ket{11}\big)\ket{-}$ & \optlabel{D}\,$\frac12\big(\ket{00} + \ket{01} + \ket{10} - \ket{11}\big)\ket{-}$
\end{tabularx}
\end{qitem}
\qsep

\needspace{7\baselineskip}
\begin{qitem}{40}
Now the ancilla alone is measured in the computational basis. The probability of the outcome $1$, and the state of the input qubits afterwards, are
\begin{itemize}[leftmargin=2.6em, labelsep=0.4em, itemsep=2pt, topsep=3pt]
\item[\optlabel{A}] $\frac14$; the input qubits collapse to $\ket{11}$
\item[\optlabel{B}] $\frac12$; the input qubits remain in $\frac12\big(\ket{00} + \ket{01} + \ket{10} - \ket{11}\big)$
\item[\optlabel{C}] $\frac12$; the input qubits collapse to $\ket{11}$
\item[\optlabel{D}] $1$; the input qubits remain in $\frac12\big(\ket{00} + \ket{01} + \ket{10} - \ket{11}\big)$
\end{itemize}
\end{qitem}
\qsep

\newpage

\begin{sectionbanner}SECTION C \enspace$\cdot$\enspace Grover Basics \hfill Q.41--60\end{sectionbanner}

{\small \textbf{Syllabus:} the angle $\theta$, oracle and diffusion circuits, inversion about the mean, success probability, optimal iterations, several solutions, over-rotation and gate counts.\\ \textbf{Formats:} Single-correct MCQ (Q.41--54); Matching lists (Q.55--56); Assertion--Reason (Q.57); Statement I / II (Q.58); Paragraph based (Q.59--60). Each question: $+3$ / $0$ / $-1$.}
\par\medskip

\needspace{7\baselineskip}
\begin{qitem}{41}
Grover's algorithm is run on $n = 2$ qubits ($N = 4$) with one marked item. After exactly one iteration, the probability of measuring the marked item is
\par\smallskip\noindent\begin{tabularx}{\linewidth}{@{}*{4}{L}@{}}
\optlabel{A}\,$\dfrac12$ & \optlabel{B}\,$\dfrac14$ & \optlabel{C}\,$\dfrac34$ & \optlabel{D}\,$1$
\end{tabularx}
\end{qitem}
\qsep

\needspace{7\baselineskip}
\begin{qitem}{42}
For $N = 8$ with one marked item, the probability of measuring the marked item after one iteration is
\par\smallskip\noindent\begin{tabularx}{\linewidth}{@{}*{4}{L}@{}}
\optlabel{A}\,$1$ & \optlabel{B}\,$\dfrac{25}{32}$ & \optlabel{C}\,$\dfrac58$ & \optlabel{D}\,$\dfrac{9}{16}$
\end{tabularx}
\end{qitem}
\qsep

\needspace{7\baselineskip}
\begin{qitem}{43}
For $n = 4$ qubits with one marked item ($\theta \approx 0.2527$ rad), the optimal number of iterations $k_{\rm opt}$ and the resulting success probability are
\par\smallskip\noindent\begin{tabularx}{\linewidth}{@{}*{2}{L}@{}}
\optlabel{A}\,$k_{\rm opt} = 2$ and $P \approx 0.908$ & \optlabel{B}\,$k_{\rm opt} = 4$ and $P \approx 0.582$ \\[3pt]
\optlabel{C}\,$k_{\rm opt} = 3$ and $P = 1$ & \optlabel{D}\,$k_{\rm opt} = 3$ and $P \approx 0.961$
\end{tabularx}
\end{qitem}
\qsep

\needspace{7\baselineskip}
\begin{qitem}{44}
The accompanying notebook searches $2^{15} = 32\,768$ items for one marked item. The optimal number of iterations is
\par\smallskip\noindent\begin{tabularx}{\linewidth}{@{}*{4}{L}@{}}
\optlabel{A}\,$181$ & \optlabel{B}\,$142$ & \optlabel{C}\,$16\,384$ & \optlabel{D}\,$128$
\end{tabularx}
\end{qitem}
\qsep

\needspace{7\baselineskip}
\begin{qitem}{45}
An oracle for the marked string $\ket{0101}$ is built as: $X$ gates, then the multi-controlled $Z$ on all four qubits, then the same $X$ gates. The $X$ gates act on
\par\smallskip\noindent\begin{tabularx}{\linewidth}{@{}*{4}{L}@{}}
\optlabel{A}\,qubits 2 and 4 & \optlabel{B}\,no qubit & \optlabel{C}\,qubits 1 and 3 & \optlabel{D}\,all four qubits
\end{tabularx}
\end{qitem}
\qsep

\needspace{7\baselineskip}
\begin{qitem}{46}
Use the circuit in the conventions box and count each multi-controlled $Z$ as one gate. The number of gates in \emph{one} Grover iteration for $n = 6$, $x^\ast = 010011$ is
\par\smallskip\noindent\begin{tabularx}{\linewidth}{@{}*{4}{L}@{}}
\optlabel{A}\,$32$ & \optlabel{B}\,$30$ & \optlabel{C}\,$26$ & \optlabel{D}\,$38$
\end{tabularx}
\end{qitem}
\qsep

\needspace{7\baselineskip}
\begin{qitem}{47}
Let $C$ be the multi-controlled $Z$ on $n$ qubits. The circuit $H^{\otimes n}X^{\otimes n}\,C\,X^{\otimes n}H^{\otimes n}$ implements
\par\smallskip\noindent\begin{tabularx}{\linewidth}{@{}*{2}{L}@{}}
\optlabel{A}\,$2\ket{s}\!\bra{s} - I = D$ & \optlabel{B}\,$I - 2\ket{s}\!\bra{s} = -D$ \\[3pt]
\optlabel{C}\,$I - 2\ket{0\dots0}\!\bra{0\dots0}$ & \optlabel{D}\,$I - 2\ket{1\dots1}\!\bra{1\dots1}$
\end{tabularx}
\end{qitem}
\qsep

\needspace{7\baselineskip}
\begin{qitem}{48}
For $N = 16$ with one marked item, the amplitude of each \emph{unmarked} item after one iteration is
\par\smallskip\noindent\begin{tabularx}{\linewidth}{@{}*{4}{L}@{}}
\optlabel{A}\,$\dfrac14$ & \optlabel{B}\,$\dfrac{11}{16}$ & \optlabel{C}\,$\dfrac{3}{16}$ & \optlabel{D}\,$\dfrac{7}{32}$
\end{tabularx}
\end{qitem}
\qsep

\needspace{7\baselineskip}
\begin{qitem}{49}
Let $N = 16$ with $M = 4$ marked items. After one iteration the probability of measuring a marked item is
\par\smallskip\noindent\begin{tabularx}{\linewidth}{@{}*{4}{L}@{}}
\optlabel{A}\,$\dfrac34$ & \optlabel{B}\,$\dfrac{25}{32}$ & \optlabel{C}\,$1$ & \optlabel{D}\,$\dfrac14$
\end{tabularx}
\end{qitem}
\qsep

\needspace{7\baselineskip}
\begin{qitem}{50}
Let $N = 8$ with $M = 4$ marked items. The probability of a marked outcome after $k$ iterations is
\par\smallskip\noindent\begin{tabularx}{\linewidth}{@{}*{2}{L}@{}}
\optlabel{A}\,$1$ for even $k$ and $0$ for odd $k$ & \optlabel{B}\,$0$ for $k = 1$ \\[3pt]
\optlabel{C}\,$1$ for $k = 1$ & \optlabel{D}\,$\frac12$ for every $k$
\end{tabularx}
\end{qitem}
\qsep

\needspace{7\baselineskip}
\begin{qitem}{51}
Let $N = 8$ with $M = 6$ marked items. The probability of a marked outcome after one iteration is
\par\smallskip\noindent\begin{tabularx}{\linewidth}{@{}*{4}{L}@{}}
\optlabel{A}\,$\dfrac34$ & \optlabel{B}\,$\dfrac14$ & \optlabel{C}\,$1$ & \optlabel{D}\,$0$
\end{tabularx}
\end{qitem}
\qsep

\needspace{7\baselineskip}
\begin{qitem}{52}
On $2^{15}$ items with one marked item, $k_{\rm opt} = 142$. If $k = 284$ iterations are used instead, the probability of the marked item is about
\par\smallskip\noindent\begin{tabularx}{\linewidth}{@{}*{4}{L}@{}}
\optlabel{A}\,$3\times10^{-6}$ & \optlabel{B}\,$\dfrac12$ & \optlabel{C}\,$\dfrac{1}{32768}$ & \optlabel{D}\,$1$
\end{tabularx}
\end{qitem}
\qsep

\needspace{7\baselineskip}
\begin{qitem}{53}
For $n = 20$ qubits with one marked item, Grover needs $804$ iterations. A classical random search needs on average about $2^{19}$ queries. The classical cost divided by the quantum cost is about
\par\smallskip\noindent\begin{tabularx}{\linewidth}{@{}*{4}{L}@{}}
\optlabel{A}\,$650$ & \optlabel{B}\,$2$ & \optlabel{C}\,$1000$ & \optlabel{D}\,$20$
\end{tabularx}
\end{qitem}
\qsep

\needspace{7\baselineskip}
\begin{qitem}{54}
A Grover circuit with success probability $P = 0.9613$ is run for $2000$ shots. The expected marked count and its binomial standard deviation are about
\par\smallskip\noindent\begin{tabularx}{\linewidth}{@{}*{4}{L}@{}}
\optlabel{A}\,$1923 \pm 44$ & \optlabel{B}\,$1923 \pm 9$ & \optlabel{C}\,$2000 \pm 0$ & \optlabel{D}\,$1923 \pm 77$
\end{tabularx}
\end{qitem}
\qsep

\needspace{7\baselineskip}
\begin{qitem}{55}
Match the number of qubits (one marked item) in List-I with $k_{\rm opt}$ in List-II. Use $k_{\rm opt} = \operatorname{round}\big(\frac{\pi}{4\theta} - \frac12\big)$.
\par\smallskip\noindent{\renewcommand{\arraystretch}{1.35}%
\begin{tabularx}{\linewidth}{@{}c>{\raggedright\arraybackslash}p{0.33\linewidth}@{\hspace{1.2em}}cL@{}}
\toprule
\multicolumn{2}{@{}l}{\textbf{List-I}} & \multicolumn{2}{l@{}}{\textbf{List-II}} \\ \midrule
\textbf{(P)} & $n = 2$ & \textbf{(1)} & $12$ \\
\textbf{(Q)} & $n = 6$ & \textbf{(2)} & $25$ \\
\textbf{(R)} & $n = 8$ & \textbf{(3)} & $1$ \\
\textbf{(S)} & $n = 10$ & \textbf{(4)} & $6$ \\
\bottomrule
\end{tabularx}}
\par\smallskip\noindent\begin{tabularx}{\linewidth}{@{}*{2}{L}@{}}
\optlabel{A}\,P$\to$3, Q$\to$4, R$\to$1, S$\to$2 & \optlabel{B}\,P$\to$3, Q$\to$1, R$\to$4, S$\to$2 \\[3pt]
\optlabel{C}\,P$\to$3, Q$\to$4, R$\to$2, S$\to$1 & \optlabel{D}\,P$\to$4, Q$\to$3, R$\to$1, S$\to$2
\end{tabularx}
\end{qitem}
\qsep

\needspace{7\baselineskip}
\begin{qitem}{56}
Match each part of the Grover circuit in List-I with its effect in List-II.
\par\smallskip\noindent{\renewcommand{\arraystretch}{1.35}%
\begin{tabularx}{\linewidth}{@{}c>{\raggedright\arraybackslash}p{0.33\linewidth}@{\hspace{1.2em}}cL@{}}
\toprule
\multicolumn{2}{@{}l}{\textbf{List-I}} & \multicolumn{2}{l@{}}{\textbf{List-II}} \\ \midrule
\textbf{(P)} & $H^{\otimes n}$ on $\ket{0\dots0}$ & \textbf{(1)} & returns one bit string with its Born probability \\
\textbf{(Q)} & $X$ gates on the $0$ bits of $x^\ast$, multi-controlled $Z$, the same $X$ gates & \textbf{(2)} & creates the uniform superposition $\ket{s}$ \\
\textbf{(R)} & $H^{\otimes n}X^{\otimes n}$, multi-controlled $Z$, $X^{\otimes n}H^{\otimes n}$ & \textbf{(3)} & reflects about $\ket{s}$, up to a global phase \\
\textbf{(S)} & measuring all qubits & \textbf{(4)} & flips the sign of $\ket{x^\ast}$ only \\
\bottomrule
\end{tabularx}}
\par\smallskip\noindent\begin{tabularx}{\linewidth}{@{}*{2}{L}@{}}
\optlabel{A}\,P$\to$2, Q$\to$4, R$\to$3, S$\to$1 & \optlabel{B}\,P$\to$2, Q$\to$4, R$\to$1, S$\to$3 \\[3pt]
\optlabel{C}\,P$\to$4, Q$\to$2, R$\to$3, S$\to$1 & \optlabel{D}\,P$\to$2, Q$\to$3, R$\to$4, S$\to$1
\end{tabularx}
\end{qitem}
\qsep

\needspace{7\baselineskip}
\begin{qitem}{57}
\textbf{Assertion (A):} The success probability of Grover's algorithm does not depend on which single item is marked.\\[1mm]
\textbf{Reason (R):} The oracle is a diagonal matrix in the computational basis.
\begin{itemize}[leftmargin=2.6em, labelsep=0.4em, itemsep=2pt, topsep=3pt]
\item[\optlabel{A}] Both Assertion (A) and Reason (R) are true, and R is the correct explanation of A
\item[\optlabel{B}] Both A and R are true, but R is \emph{not} the correct explanation of A
\item[\optlabel{C}] A is true, but R is false
\item[\optlabel{D}] A is false, but R is true
\end{itemize}
\end{qitem}
\qsep

\needspace{7\baselineskip}
\begin{qitem}{58}
\textbf{Statement I:} In Grover's algorithm, applying more iterations never lowers the success probability.\\[1mm]
\textbf{Statement II:} For one marked item, the optimal number of iterations grows as $\sqrt N$.
\begin{itemize}[leftmargin=2.6em, labelsep=0.4em, itemsep=2pt, topsep=3pt]
\item[\optlabel{A}] Both Statement I and Statement II are true
\item[\optlabel{B}] Both Statement I and Statement II are false
\item[\optlabel{C}] Statement I is true, but Statement II is false
\item[\optlabel{D}] Statement I is false, but Statement II is true
\end{itemize}
\end{qitem}
\qsep

\needspace{20\baselineskip}
\begin{passagebox}{Paragraph for Questions 59 and 60}
Grover's algorithm runs on $n = 3$ qubits with the single marked item $x^\ast = \ket{101}$. It starts from the uniform state $\ket{s}$, in which every
amplitude is $a = 1/\sqrt8$. One iteration is the oracle $O$ followed by the diffusion $D$. Trace the amplitudes by hand with inversion about the mean.
\end{passagebox}
\needspace{7\baselineskip}
\begin{qitem}{59}
After the first iteration the amplitudes are
\par\smallskip\noindent\begin{tabularx}{\linewidth}{@{}*{2}{L}@{}}
\optlabel{A}\,marked $\frac32a$; each unmarked $\frac12a$ & \optlabel{B}\,marked $\frac74a$; each unmarked $\frac34a$ \\[3pt]
\optlabel{C}\,marked $\frac52a$; each unmarked $\frac12a$ & \optlabel{D}\,marked $\frac52a$; each unmarked $\frac34a$
\end{tabularx}
\end{qitem}
\qsep

\needspace{7\baselineskip}
\begin{qitem}{60}
After the second iteration the probability of measuring $\ket{101}$ is
\par\smallskip\noindent\begin{tabularx}{\linewidth}{@{}*{4}{L}@{}}
\optlabel{A}\,$\dfrac{25}{32}$ & \optlabel{B}\,$\dfrac{11}{16}$ & \optlabel{C}\,$\dfrac{121}{128}$ & \optlabel{D}\,$1$
\end{tabularx}
\end{qitem}
\qsep

\newpage

\begin{sectionbanner}SECTION D \enspace$\cdot$\enspace Advanced Grover \hfill Q.61--80\end{sectionbanner}

{\small \textbf{Syllabus:} eigenstructure of $G$, unknown $M$, exact and fixed-point search, amplitude amplification, quantum counting, optimality, parallelism, noise, resources and cryptographic impact.\\ \textbf{Formats:} Single-correct MCQ (Q.61--74); Matching lists (Q.75--76); Assertion--Reason (Q.77); Statement I / II (Q.78); Paragraph based (Q.79--80). Each question: $+3$ / $0$ / $-1$.}
\par\medskip

\needspace{7\baselineskip}
\begin{qitem}{61}
For $N = 1024$ with $M = 16$ marked items, the optimal number of iterations is
\par\smallskip\noindent\begin{tabularx}{\linewidth}{@{}*{4}{L}@{}}
\optlabel{A}\,$25$ & \optlabel{B}\,$6$ & \optlabel{C}\,$5$ & \optlabel{D}\,$8$
\end{tabularx}
\end{qitem}
\qsep

\needspace{7\baselineskip}
\begin{qitem}{62}
When $M$ is unknown, the algorithm of Boyer, Brassard, H{\o}yer and Tapp picks $k$ uniformly from $\{0, \dots, m-1\}$. The average success probability is at least $\frac14$ once $m \ge 1/\sin2\theta$.
For $N = 2^{20}$ and $M = 1$ this threshold on $m$ is about
\par\smallskip\noindent\begin{tabularx}{\linewidth}{@{}*{4}{L}@{}}
\optlabel{A}\,$512$ & \optlabel{B}\,$1024$ & \optlabel{C}\,$804$ & \optlabel{D}\,$2^{20}$
\end{tabularx}
\end{qitem}
\qsep

\needspace{7\baselineskip}
\begin{qitem}{63}
To obtain success probability exactly $1$, the angle is reduced from $\theta$ to $\theta'$ so that $(2k+1)\theta' = \pi/2$ with $k = \lceil\frac{\pi}{4\theta} - \frac12\rceil$. For $N = 16$, $M = 1$ this requires
\par\smallskip\noindent\begin{tabularx}{\linewidth}{@{}*{2}{L}@{}}
\optlabel{A}\,$k = 3$ and $\theta' = \pi/12$ & \optlabel{B}\,$k = 4$ and $\theta' = \pi/18$ \\[3pt]
\optlabel{C}\,$k = 2$ and $\theta' = \pi/10$ & \optlabel{D}\,$k = 3$ and $\theta' = \pi/14$
\end{tabularx}
\end{qitem}
\qsep

\needspace{7\baselineskip}
\begin{qitem}{64}
Which statement about the query complexity of unstructured search is correct?
\begin{itemize}[leftmargin=2.6em, labelsep=0.4em, itemsep=2pt, topsep=3pt]
\item[\optlabel{A}] Grover's algorithm is optimal only when $N$ is a power of $2$
\item[\optlabel{B}] Quantum algorithms need only $O(\log N)$ queries, because $n = \log_2 N$ qubits hold all $N$ items
\item[\optlabel{C}] The lower bound $\Omega(\sqrt N)$ holds only for classical algorithms
\item[\optlabel{D}] Every quantum algorithm needs $\Omega(\sqrt N)$ oracle queries, so Grover's algorithm is optimal up to a constant
\end{itemize}
\end{qitem}
\qsep

\needspace{7\baselineskip}
\begin{qitem}{65}
The search space of size $N$ is split evenly among $16$ independent quantum computers, each running Grover on its share. Compared with one computer on the whole space, the wall-clock time is divided by
\begin{itemize}[leftmargin=2.6em, labelsep=0.4em, itemsep=2pt, topsep=3pt]
\item[\optlabel{A}] $16$, and the total number of oracle queries is unchanged
\item[\optlabel{B}] $16$, and the total number of oracle queries is multiplied by $4$
\item[\optlabel{C}] $4$, and the total number of oracle queries is multiplied by $4$
\item[\optlabel{D}] $4$, and the total number of oracle queries is unchanged
\end{itemize}
\end{qitem}
\qsep

\needspace{7\baselineskip}
\begin{qitem}{66}
Amplitude amplification boosts a procedure $\mathcal{A}$ that succeeds with probability $a = 0.01$ ($\theta \approx 0.10017$). At the optimal $k$, one uses $\mathcal{A}$ or $\mathcal{A}^{-1}$ in total
\par\smallskip\noindent\begin{tabularx}{\linewidth}{@{}*{4}{L}@{}}
\optlabel{A}\,$7$ times & \optlabel{B}\,$15$ times & \optlabel{C}\,$100$ times & \optlabel{D}\,$8$ times
\end{tabularx}
\end{qitem}
\qsep

\needspace{7\baselineskip}
\begin{qitem}{67}
For $N = 16$, $M = 4$, the phases $\varphi \in [0, 1)$ of the eigenvalues $e^{2\pi i\varphi}$ of $G$ on the plane containing $\ket{s}$ are
\par\smallskip\noindent\begin{tabularx}{\linewidth}{@{}*{4}{L}@{}}
\optlabel{A}\,$\dfrac14$ and $\dfrac34$ & \optlabel{B}\,$\dfrac13$ and $\dfrac23$ & \optlabel{C}\,$\dfrac1{12}$ and $\dfrac{11}{12}$ & \optlabel{D}\,$\dfrac16$ and $\dfrac56$
\end{tabularx}
\end{qitem}
\qsep

\needspace{7\baselineskip}
\begin{qitem}{68}
A cipher with a $128$-bit key is attacked by Grover's algorithm (one correct key, cost measured in iterations). The number of iterations is about
\par\smallskip\noindent\begin{tabularx}{\linewidth}{@{}*{4}{L}@{}}
\optlabel{A}\,$\frac\pi4\cdot2^{64}$ & \optlabel{B}\,$\frac\pi4\cdot2^{127}$ & \optlabel{C}\,$2^{127}$ & \optlabel{D}\,$128$
\end{tabularx}
\end{qitem}
\qsep

\needspace{7\baselineskip}
\begin{qitem}{69}
A simple noise model multiplies the signal by $f$ per iteration: $P_k \approx f^{k}\sin^2((2k+1)\theta) + (1 - f^{k})/N$. For $n = 10$, $k = 25$, $f = 0.99$ (take $0.99^{25} \approx 0.78$), the success probability is about
\par\smallskip\noindent\begin{tabularx}{\linewidth}{@{}*{4}{L}@{}}
\optlabel{A}\,$0.99$ & \optlabel{B}\,$0.50$ & \optlabel{C}\,$0.78$ & \optlabel{D}\,$0.25$
\end{tabularx}
\end{qitem}
\qsep

\needspace{7\baselineskip}
\begin{qitem}{70}
A multi-controlled $Z$ with $m$ controls can be built with $m - 1$ clean ancillas, $2(m-1)$ Toffoli gates and one CZ. For the $n = 10$ Grover register ($m = 9$), one full iteration (oracle and diffusion) needs how many Toffoli gates?
\par\smallskip\noindent\begin{tabularx}{\linewidth}{@{}*{4}{L}@{}}
\optlabel{A}\,$32$ & \optlabel{B}\,$36$ & \optlabel{C}\,$16$ & \optlabel{D}\,$18$
\end{tabularx}
\end{qitem}
\qsep

\needspace{7\baselineskip}
\begin{qitem}{71}
On a state-vector simulator the work grows as (iterations) $\times$ (amplitudes per gate) $\sim 2^{n/2}\cdot2^{n}$. Adding $2$ qubits multiplies the total work by about
\par\smallskip\noindent\begin{tabularx}{\linewidth}{@{}*{4}{L}@{}}
\optlabel{A}\,$16$ & \optlabel{B}\,$8$ & \optlabel{C}\,$4$ & \optlabel{D}\,$2\sqrt2$
\end{tabularx}
\end{qitem}
\qsep

\needspace{7\baselineskip}
\begin{qitem}{72}
In the basis $\{\ket{r}, \ket{m}\}$ of the Grover plane, $G = \begin{psmallmatrix}\cos2\theta & -\sin2\theta\\ \sin2\theta & \cos2\theta\end{psmallmatrix}$. Eigenvectors of $G$ are
\begin{itemize}[leftmargin=2.6em, labelsep=0.4em, itemsep=2pt, topsep=3pt]
\item[\optlabel{A}] $\frac{1}{\sqrt2}\big(\ket{r} \pm \ket{m}\big)$ with eigenvalues $e^{\pm2i\theta}$
\item[\optlabel{B}] $\frac{1}{\sqrt2}\big(\ket{r} \mp i\ket{m}\big)$ with eigenvalues $\pm1$
\item[\optlabel{C}] $\frac{1}{\sqrt2}\big(\ket{r} \mp i\ket{m}\big)$ with eigenvalues $e^{\pm2i\theta}$
\item[\optlabel{D}] $\ket{r}$ and $\ket{m}$ with eigenvalues $e^{\pm2i\theta}$
\end{itemize}
\end{qitem}
\qsep

\needspace{7\baselineskip}
\begin{qitem}{73}
In Grover's fixed-point ($\pi/3$) search, one recursion step turns a procedure with failure probability $\varepsilon$ into one with failure probability $\varepsilon^3$. Starting from $\varepsilon = 0.1$, two steps give
\par\smallskip\noindent\begin{tabularx}{\linewidth}{@{}*{4}{L}@{}}
\optlabel{A}\,$10^{-3}$ & \optlabel{B}\,$10^{-6}$ & \optlabel{C}\,$10^{-27}$ & \optlabel{D}\,$10^{-9}$
\end{tabularx}
\end{qitem}
\qsep

\needspace{7\baselineskip}
\begin{qitem}{74}
Let $N = 8$, $M = 1$, and repeat a run with a fixed $k$ until the marked item is found, counting only oracle queries in the Grover iterations. Then $P_1 = \frac{25}{32}$ and $P_2 = \frac{121}{128}$. The expected numbers of queries for $k = 1$ and $k = 2$ are about
\par\smallskip\noindent\begin{tabularx}{\linewidth}{@{}*{4}{L}@{}}
\optlabel{A}\,$1$ and $2$ & \optlabel{B}\,$2.12$ and $1.28$ & \optlabel{C}\,$1.28$ and $2.12$ & \optlabel{D}\,$0.78$ and $1.89$
\end{tabularx}
\end{qitem}
\qsep

\needspace{7\baselineskip}
\begin{qitem}{75}
Match each task in List-I with its query complexity in List-II.
\par\smallskip\noindent{\renewcommand{\arraystretch}{1.35}%
\begin{tabularx}{\linewidth}{@{}c>{\raggedright\arraybackslash}p{0.33\linewidth}@{\hspace{1.2em}}cL@{}}
\toprule
\multicolumn{2}{@{}l}{\textbf{List-I}} & \multicolumn{2}{l@{}}{\textbf{List-II}} \\ \midrule
\textbf{(P)} & Unstructured search, classical & \textbf{(1)} & $\Theta(N^{1/3})$ \\
\textbf{(Q)} & Unstructured search, quantum & \textbf{(2)} & $\Theta(N^{2/3})$ \\
\textbf{(R)} & Collision finding in a 2-to-1 function, quantum & \textbf{(3)} & $\Theta(N)$ \\
\textbf{(S)} & Element distinctness, quantum & \textbf{(4)} & $\Theta(N^{1/2})$ \\
\bottomrule
\end{tabularx}}
\par\smallskip\noindent\begin{tabularx}{\linewidth}{@{}*{2}{L}@{}}
\optlabel{A}\,P$\to$3, Q$\to$4, R$\to$1, S$\to$2 & \optlabel{B}\,P$\to$4, Q$\to$3, R$\to$1, S$\to$2 \\[3pt]
\optlabel{C}\,P$\to$3, Q$\to$1, R$\to$4, S$\to$2 & \optlabel{D}\,P$\to$3, Q$\to$4, R$\to$2, S$\to$1
\end{tabularx}
\end{qitem}
\qsep

\needspace{7\baselineskip}
\begin{qitem}{76}
Match each method in List-I with its main idea in List-II.
\par\smallskip\noindent{\renewcommand{\arraystretch}{1.35}%
\begin{tabularx}{\linewidth}{@{}c>{\raggedright\arraybackslash}p{0.33\linewidth}@{\hspace{1.2em}}cL@{}}
\toprule
\multicolumn{2}{@{}l}{\textbf{List-I}} & \multicolumn{2}{l@{}}{\textbf{List-II}} \\ \midrule
\textbf{(P)} & Amplitude amplification & \textbf{(1)} & uses randomly chosen $k$ from a growing range when $M$ is unknown \\
\textbf{(Q)} & Quantum counting & \textbf{(2)} & changes the reflection phases so that extra steps cannot over-rotate \\
\textbf{(R)} & BBHT search & \textbf{(3)} & replaces $H^{\otimes n}$ by any state-preparation procedure $\mathcal{A}$ \\
\textbf{(S)} & Fixed-point search & \textbf{(4)} & applies phase estimation to $G$ to find $M$ \\
\bottomrule
\end{tabularx}}
\par\smallskip\noindent\begin{tabularx}{\linewidth}{@{}*{2}{L}@{}}
\optlabel{A}\,P$\to$3, Q$\to$1, R$\to$4, S$\to$2 & \optlabel{B}\,P$\to$3, Q$\to$4, R$\to$1, S$\to$2 \\[3pt]
\optlabel{C}\,P$\to$3, Q$\to$4, R$\to$2, S$\to$1 & \optlabel{D}\,P$\to$4, Q$\to$3, R$\to$1, S$\to$2
\end{tabularx}
\end{qitem}
\qsep

\needspace{7\baselineskip}
\begin{qitem}{77}
\textbf{Assertion (A):} Grover's algorithm solves NP-complete problems in polynomial time.\\[1mm]
\textbf{Reason (R):} For a search over $2^n$ candidates, Grover's algorithm uses about $\frac\pi4 2^{n/2}$ iterations.
\begin{itemize}[leftmargin=2.6em, labelsep=0.4em, itemsep=2pt, topsep=3pt]
\item[\optlabel{A}] Both Assertion (A) and Reason (R) are true, and R is the correct explanation of A
\item[\optlabel{B}] Both A and R are true, but R is \emph{not} the correct explanation of A
\item[\optlabel{C}] A is true, but R is false
\item[\optlabel{D}] A is false, but R is true
\end{itemize}
\end{qitem}
\qsep

\needspace{7\baselineskip}
\begin{qitem}{78}
\textbf{Statement I:} The BBHT algorithm finds a marked item with an expected $O(\sqrt{N/M})$ queries without knowing $M$ in advance.\\[1mm]
\textbf{Statement II:} If exactly half of the items are marked, ordinary Grover iterations leave the success probability at $\frac12$.
\begin{itemize}[leftmargin=2.6em, labelsep=0.4em, itemsep=2pt, topsep=3pt]
\item[\optlabel{A}] Both Statement I and Statement II are true
\item[\optlabel{B}] Both Statement I and Statement II are false
\item[\optlabel{C}] Statement I is true, but Statement II is false
\item[\optlabel{D}] Statement I is false, but Statement II is true
\end{itemize}
\end{qitem}
\qsep

\needspace{20\baselineskip}
\begin{passagebox}{Paragraph for Questions 79 and 80}
\textbf{Quantum counting.} A search space of $N = 64$ items contains an unknown number $M$ of marked items; in fact $M = 16$. The Grover operator $G$ has
eigenvalues $e^{\pm 2i\theta}$ on the plane containing $\ket{s}$, with $\sin^2\theta = M/N$. Phase estimation of $G$ with $t = 3$ counting qubits is run on $\ket{s}$, which has equal weight on the two eigenvectors. The outcome $y \in \{0, \dots, 7\}$ gives the
phase estimate $\varphi \approx y/8$, where $e^{2\pi i\varphi}$ is the eigenvalue. From this one sets $\tilde\theta = \pi y/8$ and $\tilde M = N\sin^2\tilde\theta$.
\end{passagebox}
\needspace{7\baselineskip}
\begin{qitem}{79}
The two most likely outcomes are
\par\smallskip\noindent\begin{tabularx}{\linewidth}{@{}*{4}{L}@{}}
\optlabel{A}\,$y = 2$ and $y = 6$ & \optlabel{B}\,$y = 1$ and $y = 7$ & \optlabel{C}\,$y = 1$ and $y = 2$ & \optlabel{D}\,$y = 0$ and $y = 4$
\end{tabularx}
\end{qitem}
\qsep

\needspace{7\baselineskip}
\begin{qitem}{80}
If $y = 1$ is obtained, the estimate of $M$ is
\par\smallskip\noindent\begin{tabularx}{\linewidth}{@{}*{4}{L}@{}}
\optlabel{A}\,$\tilde M = 32$ & \optlabel{B}\,$\tilde M = 8$ & \optlabel{C}\,$\tilde M \approx 9.4$ & \optlabel{D}\,$\tilde M = 16$
\end{tabularx}
\end{qitem}
\qsep

\newpage

\begin{center}{\Large\bfseries\color{brand} RESPONSE SHEET}\\[2pt]
{\small Darken exactly ONE bubble per question. Leave all four bubbles empty to skip a question (0 marks).}\end{center}
\textbf{Name:} \hrulefill\quad \textbf{Roll No.:} \rule{35mm}{0.4pt}\par\medskip
\noindent\begin{minipage}[t]{0.235\linewidth}\centering
\textbf{Section A}\\[2pt]
\begin{tabular}{@{}r@{\hspace{5pt}}l@{}}
\textbf{1} & \bub{A}\,\bub{B}\,\bub{C}\,\bub{D} \\[1.6pt]
\textbf{2} & \bub{A}\,\bub{B}\,\bub{C}\,\bub{D} \\[1.6pt]
\textbf{3} & \bub{A}\,\bub{B}\,\bub{C}\,\bub{D} \\[1.6pt]
\textbf{4} & \bub{A}\,\bub{B}\,\bub{C}\,\bub{D} \\[1.6pt]
\textbf{5} & \bub{A}\,\bub{B}\,\bub{C}\,\bub{D} \\[1.6pt]
\textbf{6} & \bub{A}\,\bub{B}\,\bub{C}\,\bub{D} \\[1.6pt]
\textbf{7} & \bub{A}\,\bub{B}\,\bub{C}\,\bub{D} \\[1.6pt]
\textbf{8} & \bub{A}\,\bub{B}\,\bub{C}\,\bub{D} \\[1.6pt]
\textbf{9} & \bub{A}\,\bub{B}\,\bub{C}\,\bub{D} \\[1.6pt]
\textbf{10} & \bub{A}\,\bub{B}\,\bub{C}\,\bub{D} \\[1.6pt]
\textbf{11} & \bub{A}\,\bub{B}\,\bub{C}\,\bub{D} \\[1.6pt]
\textbf{12} & \bub{A}\,\bub{B}\,\bub{C}\,\bub{D} \\[1.6pt]
\textbf{13} & \bub{A}\,\bub{B}\,\bub{C}\,\bub{D} \\[1.6pt]
\textbf{14} & \bub{A}\,\bub{B}\,\bub{C}\,\bub{D} \\[1.6pt]
\textbf{15} & \bub{A}\,\bub{B}\,\bub{C}\,\bub{D} \\[1.6pt]
\textbf{16} & \bub{A}\,\bub{B}\,\bub{C}\,\bub{D} \\[1.6pt]
\textbf{17} & \bub{A}\,\bub{B}\,\bub{C}\,\bub{D} \\[1.6pt]
\textbf{18} & \bub{A}\,\bub{B}\,\bub{C}\,\bub{D} \\[1.6pt]
\textbf{19} & \bub{A}\,\bub{B}\,\bub{C}\,\bub{D} \\[1.6pt]
\textbf{20} & \bub{A}\,\bub{B}\,\bub{C}\,\bub{D} \\[1.6pt]
\end{tabular}
\end{minipage}\hfill\begin{minipage}[t]{0.235\linewidth}\centering
\textbf{Section B}\\[2pt]
\begin{tabular}{@{}r@{\hspace{5pt}}l@{}}
\textbf{21} & \bub{A}\,\bub{B}\,\bub{C}\,\bub{D} \\[1.6pt]
\textbf{22} & \bub{A}\,\bub{B}\,\bub{C}\,\bub{D} \\[1.6pt]
\textbf{23} & \bub{A}\,\bub{B}\,\bub{C}\,\bub{D} \\[1.6pt]
\textbf{24} & \bub{A}\,\bub{B}\,\bub{C}\,\bub{D} \\[1.6pt]
\textbf{25} & \bub{A}\,\bub{B}\,\bub{C}\,\bub{D} \\[1.6pt]
\textbf{26} & \bub{A}\,\bub{B}\,\bub{C}\,\bub{D} \\[1.6pt]
\textbf{27} & \bub{A}\,\bub{B}\,\bub{C}\,\bub{D} \\[1.6pt]
\textbf{28} & \bub{A}\,\bub{B}\,\bub{C}\,\bub{D} \\[1.6pt]
\textbf{29} & \bub{A}\,\bub{B}\,\bub{C}\,\bub{D} \\[1.6pt]
\textbf{30} & \bub{A}\,\bub{B}\,\bub{C}\,\bub{D} \\[1.6pt]
\textbf{31} & \bub{A}\,\bub{B}\,\bub{C}\,\bub{D} \\[1.6pt]
\textbf{32} & \bub{A}\,\bub{B}\,\bub{C}\,\bub{D} \\[1.6pt]
\textbf{33} & \bub{A}\,\bub{B}\,\bub{C}\,\bub{D} \\[1.6pt]
\textbf{34} & \bub{A}\,\bub{B}\,\bub{C}\,\bub{D} \\[1.6pt]
\textbf{35} & \bub{A}\,\bub{B}\,\bub{C}\,\bub{D} \\[1.6pt]
\textbf{36} & \bub{A}\,\bub{B}\,\bub{C}\,\bub{D} \\[1.6pt]
\textbf{37} & \bub{A}\,\bub{B}\,\bub{C}\,\bub{D} \\[1.6pt]
\textbf{38} & \bub{A}\,\bub{B}\,\bub{C}\,\bub{D} \\[1.6pt]
\textbf{39} & \bub{A}\,\bub{B}\,\bub{C}\,\bub{D} \\[1.6pt]
\textbf{40} & \bub{A}\,\bub{B}\,\bub{C}\,\bub{D} \\[1.6pt]
\end{tabular}
\end{minipage}\hfill\begin{minipage}[t]{0.235\linewidth}\centering
\textbf{Section C}\\[2pt]
\begin{tabular}{@{}r@{\hspace{5pt}}l@{}}
\textbf{41} & \bub{A}\,\bub{B}\,\bub{C}\,\bub{D} \\[1.6pt]
\textbf{42} & \bub{A}\,\bub{B}\,\bub{C}\,\bub{D} \\[1.6pt]
\textbf{43} & \bub{A}\,\bub{B}\,\bub{C}\,\bub{D} \\[1.6pt]
\textbf{44} & \bub{A}\,\bub{B}\,\bub{C}\,\bub{D} \\[1.6pt]
\textbf{45} & \bub{A}\,\bub{B}\,\bub{C}\,\bub{D} \\[1.6pt]
\textbf{46} & \bub{A}\,\bub{B}\,\bub{C}\,\bub{D} \\[1.6pt]
\textbf{47} & \bub{A}\,\bub{B}\,\bub{C}\,\bub{D} \\[1.6pt]
\textbf{48} & \bub{A}\,\bub{B}\,\bub{C}\,\bub{D} \\[1.6pt]
\textbf{49} & \bub{A}\,\bub{B}\,\bub{C}\,\bub{D} \\[1.6pt]
\textbf{50} & \bub{A}\,\bub{B}\,\bub{C}\,\bub{D} \\[1.6pt]
\textbf{51} & \bub{A}\,\bub{B}\,\bub{C}\,\bub{D} \\[1.6pt]
\textbf{52} & \bub{A}\,\bub{B}\,\bub{C}\,\bub{D} \\[1.6pt]
\textbf{53} & \bub{A}\,\bub{B}\,\bub{C}\,\bub{D} \\[1.6pt]
\textbf{54} & \bub{A}\,\bub{B}\,\bub{C}\,\bub{D} \\[1.6pt]
\textbf{55} & \bub{A}\,\bub{B}\,\bub{C}\,\bub{D} \\[1.6pt]
\textbf{56} & \bub{A}\,\bub{B}\,\bub{C}\,\bub{D} \\[1.6pt]
\textbf{57} & \bub{A}\,\bub{B}\,\bub{C}\,\bub{D} \\[1.6pt]
\textbf{58} & \bub{A}\,\bub{B}\,\bub{C}\,\bub{D} \\[1.6pt]
\textbf{59} & \bub{A}\,\bub{B}\,\bub{C}\,\bub{D} \\[1.6pt]
\textbf{60} & \bub{A}\,\bub{B}\,\bub{C}\,\bub{D} \\[1.6pt]
\end{tabular}
\end{minipage}\hfill\begin{minipage}[t]{0.235\linewidth}\centering
\textbf{Section D}\\[2pt]
\begin{tabular}{@{}r@{\hspace{5pt}}l@{}}
\textbf{61} & \bub{A}\,\bub{B}\,\bub{C}\,\bub{D} \\[1.6pt]
\textbf{62} & \bub{A}\,\bub{B}\,\bub{C}\,\bub{D} \\[1.6pt]
\textbf{63} & \bub{A}\,\bub{B}\,\bub{C}\,\bub{D} \\[1.6pt]
\textbf{64} & \bub{A}\,\bub{B}\,\bub{C}\,\bub{D} \\[1.6pt]
\textbf{65} & \bub{A}\,\bub{B}\,\bub{C}\,\bub{D} \\[1.6pt]
\textbf{66} & \bub{A}\,\bub{B}\,\bub{C}\,\bub{D} \\[1.6pt]
\textbf{67} & \bub{A}\,\bub{B}\,\bub{C}\,\bub{D} \\[1.6pt]
\textbf{68} & \bub{A}\,\bub{B}\,\bub{C}\,\bub{D} \\[1.6pt]
\textbf{69} & \bub{A}\,\bub{B}\,\bub{C}\,\bub{D} \\[1.6pt]
\textbf{70} & \bub{A}\,\bub{B}\,\bub{C}\,\bub{D} \\[1.6pt]
\textbf{71} & \bub{A}\,\bub{B}\,\bub{C}\,\bub{D} \\[1.6pt]
\textbf{72} & \bub{A}\,\bub{B}\,\bub{C}\,\bub{D} \\[1.6pt]
\textbf{73} & \bub{A}\,\bub{B}\,\bub{C}\,\bub{D} \\[1.6pt]
\textbf{74} & \bub{A}\,\bub{B}\,\bub{C}\,\bub{D} \\[1.6pt]
\textbf{75} & \bub{A}\,\bub{B}\,\bub{C}\,\bub{D} \\[1.6pt]
\textbf{76} & \bub{A}\,\bub{B}\,\bub{C}\,\bub{D} \\[1.6pt]
\textbf{77} & \bub{A}\,\bub{B}\,\bub{C}\,\bub{D} \\[1.6pt]
\textbf{78} & \bub{A}\,\bub{B}\,\bub{C}\,\bub{D} \\[1.6pt]
\textbf{79} & \bub{A}\,\bub{B}\,\bub{C}\,\bub{D} \\[1.6pt]
\textbf{80} & \bub{A}\,\bub{B}\,\bub{C}\,\bub{D} \\[1.6pt]
\end{tabular}
\end{minipage}
\par\bigskip
\begin{center}\textbf{Score computation} (filled in by the evaluator, using the Answer Book)\\[4pt]
\renewcommand{\arraystretch}{1.5}
\begin{tabular}{|c|l|c|c|c|c|}\hline
\textbf{Sec.} & \textbf{Topic} & \textbf{Attempted} & \textbf{Correct} $C$ & \textbf{Wrong} $W$ & \textbf{Marks} $3C - W$ \\ \hline
A & Mathematics for Grover & & & & \\ \hline
B & Quantum Mechanics & & & & \\ \hline
C & Grover Basics & & & & \\ \hline
D & Advanced Grover & & & & \\ \hline
\multicolumn{5}{|r|}{\textbf{Total (out of 240)}} & \\ \hline
\end{tabular}\end{center}

\end{document}